\section{Exercises}\label{sec:02-terminology-and-coc:03-exercises:1}

\textbf{Key points.} \texttt{Prop} is proof-irrelevant, which is exactly why \texttt{∃} (landing in \texttt{Prop}) cannot extract its witness the way \texttt{Sigma} (landing in \texttt{Type}) can. Both Π- and Σ-types, plus \texttt{Prop}, assemble into the calculus of constructions underlying every Lean declaration seen so far. β-reduction is the computational engine underneath \texttt{rfl}/\texttt{\#eval}, one substitution step at a time.

\begin{enumerate}
\def\labelenumi{\arabic{enumi}.}
\tightlist
\item
  State, as a general rule rather than by example, why \texttt{Σ n : Nat, Fin n} type-checks but \texttt{Σ n : Nat, n > 0} does not, even though \texttt{n > 0} is a perfectly good proposition about \texttt{n}. Prove the rule from the signature of \texttt{Sigma}.
\item
  \texttt{∃ x, P x} and \texttt{Σ x, P x} have exactly the same shape, a witness plus a proof. Prove that a witness cannot be extracted computationally from a proof of \texttt{∃ x, P x}, and identify exactly which fact about \texttt{Prop} the proof depends on.
\item
  β-reduce \((\lambda x.\lambda y.\, y\, x)\, a\, b\) to normal form by hand, writing out each step. The untyped-λ-calculus recap in Section 1 named \(K = \lambda x.\lambda y.\, x\) (``take two arguments, return the first''). Which existing named term does \(\lambda x.\lambda y.\, y\, x\) resemble, and how does it differ?
\item
  Construct a term of type \texttt{Σ n : Nat, Fin n} other than the \texttt{⟨3, ⟨2, by decide⟩⟩} example given in the text. Then, in a sentence or two, explain why \texttt{Σ n : Nat, n > 0} fails to type-check in Lean at all (hint, check what \emph{sort} \texttt{n > 0} lives in, and compare to the signature of \texttt{Sigma} itself).
\item
  The \texttt{Path Q : V → V → Type} of Chapter 12 was described as ``a family of types indexed by a pair of vertices.'' Write down the Π-type expression \(\prod_{x:A} B(x)\) instantiated so that it matches the signature of \texttt{Path.append}, \texttt{\{u v w : V\} → Path Q u v → Path Q v w → Path Q u w} (treat the implicit \texttt{\{u v w : V\}} as outer Π-binders). Identify \(A\) and \(B\) explicitly at each nesting level.
\end{enumerate}

Solutions, \hyperref[sec:15-appendix-solutions:02-chapter-2:1]{Appendix, Chapter 2}.
